\documentclass[11pt]{article}

\usepackage[utf8]{inputenc}
\usepackage[T1]{fontenc}
\usepackage{amsmath, amssymb, amsthm}
\usepackage{geometry}
\usepackage{graphicx}
\usepackage{hyperref}
\usepackage{enumitem}

\geometry{margin=1in}

\title{Latent Convergence and Exception-Based Learning:\\
A Unified Framework of Predictive Coherence in Artificial and Biological Cognitive Systems}

\author{Anonymous Author}

\date{}

\begin{document}

\maketitle

\begin{abstract}
Recent work has shown that artificial intelligence models trained on different modalities converge toward geometrically equivalent latent structures, a phenomenon described as the Platonic Representation Hypothesis. However, this convergence alone does not explain learning dynamics nor functional divergence across systems.

In this paper, we propose a unified framework integrating Predictive Coherence (CPEA) and Exception-Based Learning (TAE). We hypothesize that latent convergence arises from a global coherence attractor, while differentiation emerges from the handling of exceptional events.

We introduce a mathematical formalization based on latent manifolds, incoherence metrics, and exception-driven deformation dynamics. Additionally, we outline an experimental program involving embedding analysis, continual learning, and EEG correlation.

Our results suggest that intelligence—both artificial and biological—does not depend primarily on architecture, but on the interplay between structural convergence and adaptive response to exceptions.
\end{abstract}

\section{Introduction}

The scaling of artificial intelligence models has revealed an unexpected phenomenon: architectures trained on different domains converge toward similar representational structures. This challenges the classical assumption that representations are strongly dependent on data modality and model design.

The Platonic Representation Hypothesis proposes that such convergence reflects a shared underlying structure of reality. Under this view, models do not merely learn correlations but reconstruct a latent geometry common across domains.

However, this perspective has a key limitation: it describes convergence but does not explain learning dynamics or functional divergence.

In this work, we extend this hypothesis by integrating two frameworks:

\begin{itemize}
    \item Predictive Coherence (CPEA)
    \item Exception-Based Learning (TAE)
\end{itemize}

We propose the following thesis:

\begin{quote}
Any sufficiently complex cognitive system converges toward a shared latent space minimizing global incoherence, while its evolution is driven by the handling of exceptions.
\end{quote}

\section{Related Work}

\subsection{Latent Representations in AI}

Modern AI models construct high-dimensional embeddings capturing semantic relationships. Recent studies show structural similarities across models trained on different modalities.

\subsection{Predictive Processing}

Predictive processing frameworks describe cognition as continuous minimization of prediction error, widely applied in neuroscience and machine learning.

\subsection{Neuroscience of Surprise}

Experimental evidence indicates that unexpected stimuli trigger plasticity and attention mechanisms, playing a central role in learning.

\section{Theoretical Framework}

\subsection{Latent Space as a Manifold}

Let $\mathcal{M}$ be a differentiable manifold representing the latent space of a cognitive system. Each point $x \in \mathcal{M}$ corresponds to an internal representation.

We define a metric:

\begin{equation}
d(x, y)
\end{equation}

interpreted as structural incoherence between representations.

\subsection{Predictive Coherence (CPEA)}

We define the coherence functional:

\begin{equation}
\mathcal{C} = \mathbb{E}[\|x - \hat{x}\|^2]
\end{equation}

where $\hat{x}$ is the predicted representation.

Minimizing $\mathcal{C}$ leads to stabilization and convergence of the latent geometry.

\subsection{Exception-Based Learning (TAE)}

We define exceptional events as:

\begin{equation}
E_t = \{x_t : \|x_t - \hat{x}_t\| > \epsilon\}
\end{equation}

These events represent persistent structural deviations.

The update dynamics are given by:

\begin{equation}
\Delta \mathcal{M} \propto E_t
\end{equation}

Thus, the latent manifold deforms locally in response to exceptions.

\subsection{Theoretical Implication}

From these definitions:

\begin{itemize}
    \item Global convergence is driven by coherence minimization
    \item Functional differentiation emerges from exception dynamics
\end{itemize}

\section{Hypotheses}

\begin{itemize}
    \item H1: Distinct models converge toward isomorphic latent structures
    \item H2: Convergence corresponds to minimization of predictive incoherence
    \item H3: Model differences emerge in regions of high exception density
    \item H4: EEG signals exhibit analogous dynamics to latent deformation
\end{itemize}

\section{Experimental Design}

\subsection{Latent Convergence Analysis}

We compare embeddings from vision and language models using:

\begin{itemize}
    \item Centered Kernel Alignment (CKA)
    \item Procrustes analysis
    \item Representational Similarity Analysis (RSA)
\end{itemize}

\subsection{Exception Injection}

We introduce controlled anomalies into datasets and measure:

\begin{itemize}
    \item Local deformation of embeddings
    \item Persistence of structural changes
\end{itemize}

\subsection{CPEA Architecture}

We implement a model with coherence-based loss and track stability over time.

\subsection{EEG Integration}

We analyze EEG signals under expected and unexpected stimuli and compare them to artificial embeddings.

\section{Implementation}

\subsection{Technology Stack}

\begin{itemize}
    \item PyTorch
    \item Avalanche (continual learning)
    \item snnTorch (optional)
\end{itemize}

\subsection{Modular Architecture}

\begin{enumerate}
    \item Encoder (embedding generation)
    \item Predictor (estimation of $\hat{x}$)
    \item Metric (incoherence computation)
    \item Exception detector
    \item Manifold updater
\end{enumerate}

\section{Expected Results}

We expect:

\begin{itemize}
    \item Structural convergence across models
    \item Localization of exception-driven deformation
    \item Partial correspondence with EEG dynamics
\end{itemize}

\section{Discussion}

\subsection{Artificial Intelligence}

The results suggest that architecture is secondary to learning dynamics.

\subsection{Neuroscience}

The brain can be interpreted as a system balancing coherence and exception adaptation.

\subsection{Epistemology}

Reality emerges as a stable solution in representation space rather than as direct input.

\section{Limitations}

\begin{itemize}
    \item Difficulty in defining true exceptions
    \item Noise in biological signals
    \item Metric dependence
\end{itemize}

\section{Future Work}

\begin{itemize}
    \item Multimodal integration
    \item Real-time adaptive systems
    \item EEG-AGI closed-loop systems
\end{itemize}

\section{Conclusion}

We propose a unified framework combining latent convergence and exception-based learning, showing that:

\begin{quote}
Intelligence emerges from the interaction between structural stability and adaptive response to exception.
\end{quote}

This perspective shifts the focus from architecture design to learning dynamics.

\section*{References}

\begin{itemize}
    \item Representation convergence studies (arXiv, 2024)
    \item Predictive processing literature (Friston et al.)
    \item Neuroscience of surprise and plasticity
    \item Continual learning frameworks (Avalanche)
\end{itemize}

\end{document}
